\begin{enumerate}
	\item On constate que Tahiti est à peu près au milieu entre l'équateur et le parallèle 30\degres{}S, et un peu plus à l'est que le méridien 150\degres{}O.

	Le décalage entre Tahiti et le méridien 150\degres{}O est comparable au décalage entre Los Angeles et le méridien 120\degres{}O, donc on peut dire que les coordonnées géographiques approximatives de Tahiti sont $(148\degres{}\mathrm{O} ; 15\degres{}\mathrm{S})$.

	\emph{Remarque :} L'île de Tahiti est en réalité aux coordonnées $(149\degres{}\mathrm{O} ; 17\degres{}\mathrm{S})$, mais la lecture sur le graphique ne permet pas d'avoir une précision au degré près, il y a donc une gamme de valeurs acceptables, comme réponse.

	\item Le trajet a duré 22 h 10, temps d'attente compris.

	Si on enlève 2 h 20 de temps d'attente, cela fait :

	$22\text{ h }10 - 2\text{ h }20 = 22\text{ h }10 - 2 \text{ h }10 - 0\text{ h }10 = 20\text{ h} - 0\text{ h }10 = 19\text{ h }50$.

	La durée combinée des deux vols est donc de 19 h 50 min.

	\item \begin{enumerate}
		\item Le diamètre de la médaille étant de \np[cm]{8,5}, son rayon est donc la moitié de cette longueur : \np[cm]{4,25}.

		En appliquant la formule rappelée au document 3 :

		$\mathcal{V}_{\text{médaille}} = \pi \times 4,25^2 \times 0,92 = \np{16,6175}\pi \approx52,205$.

		Comme toutes les longueurs étaient en centimètres, on a bien, au dixième près un volume de \np[cm^3]{52,2}.

		\item La masse d'argent sera donc proportionnelle à son volume. En faisant le calcul avec le volume arrondi au dixième obtenu précédemment :\quad
		$10,5 \times 52,2 = 548,1$.

		D'après ce calcul, la médaille contient donc environ \np[g]{548} d'argent, au gramme près.

		\emph{Remarque :} En réalité, les médailles d'argent des jeux de Paris 2024 sont faites d'argent (pureté 925/1000), pour environ 507 g, et de 18 g de fer issu des rénovations de la Tour Eiffel, pour un total de 525 g.

		Notre calcul donne un résultat supérieur, car il ne prend pas en compte les motifs en creux de la médaille, qui font qu'il y a un peu moins d'argent que le cylindre plein.


	\end{enumerate}
\end{enumerate}
